\subsection{مبانی نظری و فرمول‌بندی ریاضی}
فرآیند فشرده‌سازی ماژول \lr{optimal\_subspace\_projection} بر پایه چهار گام ریاضی منسجم بنا شده است:
\begin{enumerate}
    \item \textbf{محاسبه ماتریس کوواریانس فعال‌سازی و منظم‌سازی:}
    فرض کنید تانسور فعال‌سازی ورودی به یک لایه خطی با $N$ نمونه توکن به صورت $\mathbf{X} \in \mathbb{R}^{N \times d_{\text{in}}}$ باشد. ماتریس کوواریانس به همراه مقدار منظم‌سازی تیخونوف $\epsilon$ جهت تضمین مثبت معین بودن محاسبه می‌شود:
    \begin{equation}
    \mathbf{C} = \frac{1}{N} \mathbf{X}^\top \mathbf{X} + \epsilon \mathbf{I}
    \end{equation}
    که در آن مقدار منظم‌سازی متناسب با میانگین مقادیر قطری $\epsilon = \max(10^{-4} \cdot \operatorname{mean}(\operatorname{diag}(\mathbf{C})), 10^{-5})$ تنظیم می‌گردد.

    \item \textbf{پیش‌شرط‌سازی چولسکی (\lr{Cholesky Preconditioning}):}
    ماتریس کوواریانس به مولفه مثلثی پایینی تجزیه می‌شود:
    \begin{equation}
    \mathbf{C} = \mathbf{L} \mathbf{L}^\top
    \end{equation}
    سپس ماتریس وزن متراکم اولیه $\mathbf{W} \in \mathbb{R}^{d_{\text{out}} \times d_{\text{in}}}$ با اعمال فاکتور چولسکی تغییر مقیاس می‌یابد تا تاثیر ویژگی‌های پرانرژی در فضای بازنمایی نمایان گردد:
    \begin{equation}
    \mathbf{M} = \mathbf{W} \mathbf{L}
    \end{equation}

    \item \textbf{تکرارهای توانی و پروجکشن زیرفضای تصادفی (\lr{Randomized Subspace Iteration}):}
    برای ماتریس $\mathbf{M}$، یک ماتریس گاوسی تصادفی $\mathbf{\Omega} \in \mathbb{R}^{d_{\text{in}} \times r}$ با رتبه هدف $r = \lfloor \min(d_{\text{in}}, d_{\text{out}}) \cdot \rho \rfloor$ نمونه‌برداری می‌شود. سپس برای تقویت طیف تکین، تکرارهای توانی صورت می‌پذیرد:
    \begin{equation}
    \mathbf{Y} = \mathbf{M} \mathbf{\Omega}, \quad \mathbf{Y} \leftarrow \mathbf{M} (\mathbf{M}^\top \mathbf{Y}) \quad (\text{تکرار به تعداد } p \text{ بار})
    \end{equation}
    سپس پایه یکامتعامد $\mathbf{Q}$ از طریق تجزیه \lr{QR} استخراج می‌شود:
    \begin{equation}
    \mathbf{Q}, \mathbf{R} = \operatorname{qr}(\mathbf{Y})
    \end{equation}

    \item \textbf{تجزیه کم‌مرتبه و معکوس‌سازی پیش‌شرط:}
    پروجکشن کوچک‌شده $\mathbf{B} = \mathbf{Q}^\top \mathbf{M} \in \mathbb{R}^{r \times d_{\text{in}}}$ تحت \lr{SVD} فشرده قرار می‌گیرد:
    \begin{equation}
    \mathbf{B} = \mathbf{U} \mathbf{\Sigma} \mathbf{V}^\top
    \end{equation}
    در نهایت با گزینش $r$ مولفه اول و خنثی‌سازی پیش‌شرط چولسکی با $\mathbf{L}^{-1}$، بازسازی ماتریس فشرده به دست می‌آید:
    \begin{equation}
    \mathbf{W}_{\text{osp}} = (\mathbf{Q} \mathbf{U}_r \mathbf{\Sigma}_r \mathbf{V}_r^\top) \mathbf{L}^{-1}
    \end{equation}
\end{enumerate}
